\documentclass[a4paper]{report}

\usepackage{amsmath}
\usepackage{graphicx}

\usepackage{anyfontsize}

\usepackage{fancyhdr}
\pagestyle{fancy}

\usepackage{titlesec}
\titleformat{\chapter}[block]
{\normalfont\huge\bfseries}{\thechapter.}{1em}{\Huge}
\titlespacing*{\chapter}{0pt}{-19pt}{19pt}

\usepackage[nottoc,numbib]{tocbibind}


\renewcommand{\familydefault}{\sfdefault} % sans serif default

\usepackage[none]{hyphenat}

\setcounter{tocdepth}{1} % Show sections
\setcounter{tocdepth}{2} % + subsections
%\setcounter{tocdepth}{3} % + subsubsections
%\setcounter{tocdepth}{4} % + paragraphs
%\setcounter{tocdepth}{5} % + subparagraphs

\title{Aplicació del Material Point Method (MPM) a la deformació d'objectes}

\author{Martín Garcia, Pol}
\date{\parbox{\linewidth}{\centering%
		2019, Quadrimestre de tardor\endgraf\bigskip
		Director: Susín Sánchez, Toni\endgraf \medskip
		Especialitat Computació \endgraf
		Facultat d'Informàtica de Barcelona \endgraf UPC}}

\lhead{\includegraphics[width=5cm]{images/logo-fib.png}}
\rhead{\fontsize{5}{6}\selectfont{\textbf{MPM a la deformació d'objectes} \\
			Lliurament 1 GEP\\}}
		
		
\begin{document}
	\pagenumbering{gobble}
	\maketitle
	\newpage
	

	\renewcommand{\contentsname}{Índex}
	
	\tableofcontents
	\newpage
	
	\pagenumbering{arabic}	
	
	\chapter{Introducció}
	Aquest treball de fi de grau ha estat centrat en la simulació de fluids, una temàtica molt àmplia que tot i que conceptualment sembla trivial, és molt complexa. Molts detalls i decisions s'han de tenir en compte a l'hora d'implementar un simulador d'aquestes característiques, com la representació del fluid, si aquest és compressible, elàstic, interactiu amb col·lisions, etc. \newline
	Per a poder presentar el treball correctament, aquest document també serveix com a introducció a la simulació de fluids sense assumir coneixements previs en aquest aspecte, i des del punt de vista d'un enginyer informàtic.\newline
	I finalment també es busca mostrar com materials sòlids-elàstics es poden representar i simular com a fluids.\newline

	\section{Context}
	La simulació de fluids, o dinàmiques de fluids computacionals (CFD's), és una disciplina que, mitjançant tècniques d'anàlisi numèrica, busca resoldre problemes que impliquen d'alguna manera un tipus de fluctuació o interacció amb fluids. \newline
	La base de qualsevol simulador de fluids són les equacions de Navier-Stokes, que descriuen la mecànica d'aquests; però necessitem molt més que les equacions per a poder implementar aquest programa, doncs també són necessaris coneixements de software i hardware per a una implementació eficient, tècniques de gràfics per computador per a una adequada visualització, i coneixements matemàtics d'àlgebra lineal que ja s'han vist durant l'educació oferida per la Facultat d'Informàtica de Barcelona, sobretot en la branca de computació. Tot i això, l'abast del projecte ens obliga a sortir de l'àmbit de coneixements d'un enginyer informàtic, i integrar múltiples conceptes matemàtics de diferents caires, que seran explicats degudament. \newline
	En el nostre cas les matemàtiques ens ajudaran a definir el comportament i les limitacions del nostre simulador, que a partir d'ara anomenarem també \textit{solver}, per a poder configurar-lo amb coneixement a posteriori, i sempre ser conscients de l'estat del sistema que s'està processant. \newline
	
	\subsection{Conceptes bàsics}
	Abans d'entrar en detalls, hi ha uns coneixements bàsics que s'han de tenir en compte degut a l'àmplia projecció dels simuladors de fluids.
	\subsubsection[Tipus de solvers]{Tipus de solvers}
	Podem diferenciar els tipus de solvers en 3 subgrups, que varien intrínsecament en el mateix concepte de la representació del fluid i els algorismes emprats. Aquesta representació marcarà la manera d'interactuar i tractar el fluid tant amb si mateix, com amb cossos externs.
	
	\paragraph[Graella]{Solvers en Graella} Podem representar el fluid en un instant de temps concret com a una magnitud en un punt d'una graella, identificant la quantitat de fluid en aquella posició en un moment concret; per altra banda, cada cel·la de la graella té alguna representació de direcció i magnitud de moviment del fluid, per a poder computar el desplaçament d'aquest. \newline
	A més a més, podem emmagatzemar altra informació a cada posició de la graella, o usar les dades ja guardades per a processar l'estat del sistema en un instant de temps posterior, tenint en compte que a cada cel·la només hi pot haver una quantitat màxima de fluid (volum màxim) perquè no sigui un fluid comprimible, i així provocar una dissipació d'aquest i el consegüent moviment.\newline
	En conjunt, creem un espai acotat per les dimensions de la graella on el fluid es mou de manera quantitativa a través de les diferents cel·les, d'acord amb la mateixa informació d'aquestes. \newline
	Aquesta representació es coneix com a simulador Eulerià.
	
	\paragraph[Partícules]{Solvers de Partícules} Podem representar un petit volum de fluid com a una partícula en una posició determinada en un instant de temps, de manera que sempre representa la mateixa quantitat de fluid, i a més guardem tota la informació de moviment (força, velocitat, ...) a un nivell molt concret representat per la mateixa partícula. \newline
	És important entendre que les partícules no representen el fluid a escala de molècules o àtoms, sinó que cada partícula representa una porció contínua del material o fluid, o un subconjunt del domini físic a simular.
	\newline
	Aquesta representació permet, per exemple, barrejar dos fluids mantenint alhora les seves propietats separades, o simular interaccions entre partícules o sòlids externs amb un major nivell de detall.
	També cal dir, que aquest mètode dificulta la tasca de mantenir un volum constant del fluid, i sobretot necessita calcular interaccions entre totes les partícules, còmput que resulta molt costós. \newline
	Els simuladors de partícules són coneguts com a Lagrangians.
	
	\paragraph[Híbrids]{Solvers Híbrids} Si mesclem els dos conceptes anteriors, obtenim un simulador que, podem dir, rep el millor dels dos mètodes.\newline
	Per una banda, es representa el fluid com partícules en cada instant de temps amb totes les seves propietats; i per altra banda el còmput d'interaccions entre aquestes el gestionem mitjançant una graella (la qual defineix l'espai del sistema) a on traspassem les característiques de les partícules en determinades zones, de manera que la computació de col·lisions i/o interaccions és molt més eficient i acotada.
	
	\subsubsection{Propietats del fluid} 
	Un fluid sempre ha de tenir representades, en alguna determinada estructura o directament en l'algorisme, un seguit de propietats que el defineixen i n'especifiquen el comportament. Cal dir, que només és necessari usar aquelles propietats que el nostre model té en compte.
	\paragraph{Volum} La quantitat de fluid que volem simular té un impacte directe en el nostre \textit{solver}, i depenent de la representació del fluid aquest concepte afectarà de manera diferent; per exemple, en simuladors de graella, massa fluid pot saturar el model; i en simuladors de partícules hem de controlar la proximitat entre partícules perquè aquesta defineix el volum.
	\paragraph{Massa} Depenent del model a implementar, ens pot interessar la massa del fluid (representada, per exemple, en densitat) per a tractar la interacció amb forces externes (col·lisions o gravetat).
	\paragraph{Elasticitat} L'elasticitat és l'habilitat d'un cos o material de recuperar la seva forma original després d'una deformació, i alhora resistir-se a aquestes forces de deformació. Es descriu amb el model i la llei de Hook, que estipula que la força de retorn és proporcional a l'extensió de la deformació.
	\paragraph{Viscositat} Molt semblant a l'elasticitat, la viscositat d'un fluid és la resistència que té a deformar-se o, informalment, a escampar-se. Es mesura en pascals per segon ($\frac{Pa}{s}$ o $\frac{kg}{m \cdot s}$) i es representa amb $\mu$ o $\eta$. És molt comú representar aquesta magnitud en relació amb la densitat $\rho$ de manera $\nu = \frac{\mu}{\rho}$ per a simplificar expressions.
	\paragraph{Plasticitat} Si l'elasticitat defineix la possible deformació d'un cos, la plasticitat marca el punt de no retorn d'una deformació, en què les forces que busquen el retorn de l'objecte a l'estat original són permanentment degradades, de manera que part de la deformació es torna permanent. La idea de modelar aquesta característica, i el pioner treball de Terzopoulos et al. \cite{Terzopoulos:1988:MID:378456.378522} el 1988, ha fet que la simulació de materials topològicament canviants sigui una popular àrea d'estudi a gràfics per computador, on es poden tractar fets com la fractura (límit de la plasticitat), i el tallament d'objectes deformables. Totes aquestes interaccions s'anomenen inelàstiques.
	\paragraph{Viscoelasticitat} Propietat dels materials que són viscosos i elàstics alhora en el moment de deformar-se. La viscositat fa que, a mesura que el temps avança, el material sofreix una tensió que li provoca una certa deformació, de manera que tota deformació elàstica sempre comporta una pèrdua d'energia a causa de la viscositat, produint una deformació plàstica.
	
	\subsubsection[MPM]{Material Point Method}
	Un dels mètodes que ara ha ressorgit per a la simulació de fluids és el conegut \textit{Material Point Method}, proposat per Sulsky \cite{Sulsky1995}, basat en partícules Lagrangianes i una graella Euleriana de rerefons. Aquest mètode és usat donat la seva alta eficiència de còmput, i la possibilitat de detallar la seva precisió amb la mida de la graella. \newline
	Aquesta tècnica intenta emmagatzemar les dades de les partícules interpolades en les cel·les més properes, du a terme els càlculs en la graella, i finalment es fa la interpolació de manera inversa i es produeix el desplaçament de la partícula. \newline
	Tot i això MPM té alguns problemes, ja que per culpa de la simplificació que es produeix a la graella no som capaços de representar velocitats discontínues, o col·lisions molt concretes amb objectes de menor resolució de la graella. 
	
	\subsection{Formulació del problema}
	Podríem considerar que hi ha dos tipus de simulacions: les que tenen com a objectiu emular la realitat, i les que busquen obtenir una animació visualment realista.\newline
	Les simulacions que imiten la realitat tendeixen a requerir moltíssima precisió i algorismes poc optimitzables, les quals acostumen a necessitar ser processades en un supercomputador per obtenir resultats adequats en un temps raonable. \newline
	Per altra banda, la simulació amb objectius només visuals (com és el nostre cas) intenta ser ràpida, visualment realista i, sobretot, configurable per a poder triar quins resultats es volen aconseguir. Encara que no s'intenti emular la realitat al cent per cent, és important imitar-la al màxim acceptant generalitzacions i truncaments per a accelerar la velocitat de computació; depenent de la situació podríem simplificar el model el suficient per arribar a tenir una simulació en temps real.\newline
	A causa del fet que es busca trobar una fidel i ràpida aproximació a la realitat, no existeix una solució analítica (ni general ni específica) per a produir aquestes simulacions. També s'ha de tenir en compte quina és la versatilitat de la simulació, perquè el mateix algorisme implementat limitarà les propietats del fluid o sòlid. Finalment també hi ha la possibilitat, i el problema, de fer interaccionar el fluid a simular amb algun cos extern (tant en moviment com immòbil) i la transmissió de forces entre aquests. \newline
	Així doncs, aquest treball pretén estudiar com plantejar solucions per a implementar les qüestions anteriors, i alhora aplicar-les per a conèixer les limitacions i habilitats de cadascuna, en un determinat sistema; i també servir com a base introductòria a aquesta temàtica tan dispersa i tècnica a qualsevol lector.
	
	\subsection{Actors implicats}
	En aquesta secció es definiran els actors implicats, de manera directa i indirecta, en el projecte.
	\paragraph{Director} Figura essencial pel correcte desenvolupament del projecte, i màxim responsable en la guia i consell al desenvolupador del treball. La seva acció és clau per determinar errors en el projecte, tant de tipus proposicional com executius. En aquest cas el director és Toni Susín Sánchez, del departament de Matemàtiques Aplicades de la UPC, i cap del \textit{Dynamic Simulation Lab} inclòs en el grup de recerca ViRVIG.
	\paragraph{Desenvolupador} La persona a càrrec de dur a terme la recerca, documentació i implementació del software requerit, així com de redactar la memòria del projecte sota la guia del director. També recau sobre aquest actor la gestió del projecte, i és en última instància la persona responsable de complir amb els terminis del treball. Aquest rol recau sobre Pol Martín Garcia, alumne de la FIB de la branca de computació.
	\paragraph{Entorn} El centre de treball ViRVIG (Visualització, Realitat Virtual i Interacció Gràfica) ha facilitat un entorn de treball per a poder usar i provar el projecte amb la potència de còmput necessària, i aprovisionat d'un clima de treball per al correcte desenvolupament del projecte.
	\paragraph{Usuaris} Aquest treball busca crear una dissertació i ampliació dels simuladors d'avui dia, i com a tal no tindrà usuaris directes. Qualsevol avanç en \textit{solvers} de tota mena, però, té un impacte parcial en la investigació i implementació de noves tècniques d'animació; com per exemple en efectes especials de pel·lícules, com a \textit{Big Hero 6}, \textit{Zootopia}, o inclús \textit{Frozen} de Disney \cite{Stomakhin} per a simular la neu.
	\paragraph{Beneficiaris} Perquè aquest projecte no crearà un producte, com s'ha dit abans, tampoc hi ha beneficiaris directes. No obstant això, altres estudiants d'enginyeria informàtica o matemàtiques poden trobar aquest projecte interessant com a base d'aprenentatge, i usar-lo com a punt de partida per a futurs projectes centrats en l'àmbit de simulació de fluids en MPM.
	
	\section{Justificació}
	Aquesta temàtica és una àmplia àrea d'estudi des de fa diverses dècades, que s'ha anat dividint en múltiples tècniques i cada branca d'estudi va aprofundint en paral·lel de la resta.
	\newline
	En aquest treball estudiarem i implementarem un \textit{solver} MPM, que neix de la generalització de \textit{Particle in Cell} (PIC) i \textit{Fluid Implicit Particle Method} (FLIP) \cite{Sulsky1995}. Hi ha hagut molt nombrosos i extensos treballs de FLIP en els darrers anys \cite{Bridson2018,Zhu2005}, mentre que MPM ha sigut realment introduït i començat a estudiar recentment perquè, al contrari que FLIP, aquest pot tractar amb fluids i sòlids comprimibles.
	\newline
	D'aquesta manera, MPM s'està començant a estudiar i a usar per a simular diferents tipus d'interaccions innovadores. Com s'ha dit abans, ressalta la simulació de neu de Disney \cite{Stomakhin} o l'interacció entre o amb múltiples materials deformables\cite{Hegemann2013}, per posar alguns exemples. Es pot trobar més informació dels treballs en MPM a \cite{Jiang2016}.
	\newline
	Just molt recentment ha sorgit una nova tècnica per a resoldre MPM, anomenada MLS-MPM (\textit{Moving Least Squares Material Point Method})\cite{hu2018mlsmpmcpic} que, generalment, busca millors aproximacions a la graella gràcies a una B-spline\footnote{Funció corba; és detallarà en futurs apartats.} calibrada amb MLS.
	\newline
	D'aquesta manera, i amb base dels nous desenvolupaments de Yuanming Hu et. al. a \cite{hu2018mlsmpmcpic} i usant com a base \cite{Hu} dels mateixos autors, estudiarem aquesta tècnica i possibles diferents interpretacions del model. És a dir: adaptarem la solució de Yuanming a unes necessitats personals i investigar com reacciona el sistema a unes noves condicions.
	
	\section{Abast}
	En aquesta secció definirem quines són les fites del projecte, i la seva estructura progressiva.
	\subsection{Objectius}
	\paragraph{Aprenentatge de nous conceptes} Donat que el treball requereix de coneixements allunyats dels adquirits en una enginyeria informàtica, cal una cerca intensiva d'informació i aprenentatge de mètodes numèrics, càlcul, i àlgebra lineal. Tots aquests conceptes són necessaris per entendre els diferents models d'objectes i simuladors a implementar. \newline
	En concret, existeixen alguns subtemes que mereixen una menció especial com a subobjectius:
	\subparagraph[Navier-Stokes]{$\cdot$ Les equacions de Navier-Stokes} Aplicant la segona llei de Newton al moviment del fluid, amb algunes assumpcions de les forces internes (anomenades estrès), s'obtenen les equacions de Navier-Stokes; conegudes per la seva complexitat, que ja de per si obren una interessant àrea d'estudi de possibles implicacions i aplicacions.
	\newline
	És important entendre aquestes equacions, i saber interpretar-les, doncs descriuran el nostre sistema i són la base de qualsevol \textit{solver}.
	\subparagraph[Càlcul vectorial]{$\cdot$ Càlcul vectorial} Fins a un cert nivell, és important i necessari conèixer diferenciació i integració de camps vectorials, doncs tractarem amb camps de partícules en aquest projecte.
	\subparagraph[Tècniques de solvers]{$\cdot$ Tècniques de solvers} Finalment, amb els coneixements teòrics necessaris, s'ha d'aprendre el funcionament intern i els models que permeten simular els diferents tipus de \textit{solvers} creats, en concret el de MLS-MPM \cite{hu2018mlsmpmcpic}.
	
	\paragraph{Implementació d'un solver bidimensional} Per a una primera versió, s'implementarà un simulador de fluids de partícules bidimensional, seguint\cite{Hu,hu2018mlsmpmcpic} i programant la base de tot el projecte a una escala menor, per a fer proves i trobar errors. La idea és usar aquesta primera fase com a indicador, i prendre decisions d'acord amb els resultats obtinguts per a millorar el resultat final, i evitar errades. \newline
	L'objectiu d'aquesta etapa és el refinament de l'entorn de treball per a una posterior fase molt més complexa, i d'aquesta manera poder evitar nombrosos possibles errors, estudiar inconvenients que puguin aparèixer i, sobretot, com a guia d'aprenentatge progressiu de manera simplificada. 
	
	\paragraph{Millora del solver a la tercera dimensió} El canvi de dimensió és un salt difícil, i essencial en aquest projecte, doncs tot simulador gràfic per a resultats realistes requereix aquesta tercera dimensió. Serà necessari redimensionar la interfície gràfica, les estructures de dades, i el mateix algorisme.
	\newline
	A partir d'aquesta fita, és quan s'inicien els tests i objectius secundaris.
	\subparagraph[Interacció viscoelàstica]{$\cdot$ Interacció entre diversos materials viscoelàstics diferents} Afegir suport per a poder mantenir en el simulador alhora diversos fluids de diferents característiques cadascun.
	\subparagraph[Sistema de físiques]{$\cdot$ Sistema de físiques} Implementar un sistema d'interacció de cossos sòlids en moviment amb el solver, per aconseguir simular col·lisions amb un, o diversos, elements externs al simulador; com per exemple una roda amb sorra.
	
	\paragraph{Altres objectius} En la mesura del possible, s'intentarà dur a terme els següents objectius, que no són necessaris per al correcte desenvolupament del projecte, però el completen i ajuden a veure com tractar i solucionar altres problemes inherents en qualsevol sistema programat.
	\subparagraph[Eficiència]{$\cdot$ Eficiència} Sobretot en el \textit{solver} tridimensional, l'eficiència i velocitat del programa és un problema, ja que aproximadament el còmput de cada fotograma serà major a un minut segons els articles de MPM \cite{hu2018mlsmpmcpic}; per tant sempre que es pugui s'ha d'implementar el simulador sent conscients de l'arquitectura i de l'eficiència del codi.
	\subparagraph[Paral·lelisme]{$\cdot$ Paral·lelisme} Aconseguir que el codi funcioni de manera paral·lela pot suposar un increment de la velocitat de computació molt important. Així doncs, qualsevol oportunitat de paral·lelitzar el codi s'haurà d'aprofitar.
	
	
	\subsection{Requisits}
	Podem separar els requisits del software a implementar en funcionals, i no funcionals. Els primers representen funcionalitats específiques que el nostre sistema haurà de complir, i el segon especifica criteris per jutjar els comportaments específics d'aquest.
	
	\subsubsection{Requisits funcionals}
	\paragraph{Simulació d'un sistema personalitzable} Amb modificació del nombre, tipus, i posició de les partícules; canvi del medi (gravetat, mida de la graella, etc). Cal valorar si han de ser paràmetres modificables entre execucions, o constants en el codi.
	\paragraph{Exportació de la simulació} Exportació dels fotogrames produïts pel solver, o bé en arxius binaris per al seu posterior tractament, o bé en imatges PNG (\textit{Portable Network Graphics}).
 	\paragraph{Simulació d'objectes viscoelàstics} Importació de models de triangles, i simular-ne el comportament en unes condicions especificades com a materials elàstics o viscoelàstics.
	
	\subsubsection{Requisits no funcionals}
	\paragraph{Velocitat de computació} La velocitat del còmput ha d'estar justificada i ser adient; ha de ser linealment dependent del nombre de partícules o a la mida de la graella.
	\paragraph{Estabilitat de la simulació} Les solucions del \textit{solver} han de ser estables, i no descontrolar-se; en cas de ser impossible evitar la inestabilitat, aquest fet estarà documentat i estudiat, marcant els límits del simulador.
	
	
	\subsection{Obstacles i riscos}
	Per evitar imprevistos en el transcurs del projecte, analitzarem els obstacles que poden sorgir durant aquest procediment i, sobretot, com minimitzar el seu impacte.
	\subsubsection{Errors en el codi}
	És inevitable cometre petits errors en la implementació d'un software, que a vegades poden passar desapercebuts, i altres vegades poden trencar totalment el programa.\newline
	Donat que un simulador no pot ser testejat amb jocs de proves, serà necessari provar manualment cada nova característica afegida, i comprovar que totes les anteriors propietats segueixen funcionant adequadament. D'aquesta manera podrem minimitzar els riscos de perpetuar un error durant les diverses versions del software.
	\subsubsection{Problemes o insuficiència de hardware}
	Els ordinadors poden tenir problemes, els discos poden fallar, la memòria RAM pot deixar de funcionar, els arxius es poden corrompre, etc. El qual suposa una pèrdua de temps important a curt termini per a arreglar el hardware espatllat i recuperar les dades perdudes. \newline
	Per evitar l'impacte d'aquest problema en la mesura de tot el possible, el projecte i la memòria d'aquest es guardaran per versions en un repositori a la xarxa, més concretament a \textit{GitHub}. A més a més, la memòria es guardarà en paral·lel al \textit{OneDrive} de Microsoft.
	\newline
	Per altra banda, durant les etapes més avançades del projecte el temps de simulació es veurà incrementat quantitativament, cosa que provocarà que per provar els canvis siguin necessàries hores. \newline
	Per a evitar perdre temps sense avançar adequadament, els tests (en forma d'execucions arbitràries del programa) es deixaran executant en una màquina diferent de la que es treballa, amb molta més potencia (com un dels ordinadors del ViRVIG) per obtenir els resultats tan aviat com sigui possible, sense entorpir el desenvolupament.
	\subsubsection{Errors de comprensió dels models matemàtics}
	A causa de la complexitat de les matemàtiques que hi ha darrere d'aquest projecte, i al fet que aquest és conceptualitza des del punt de vista informàtic, és possible errar en l'enteniment de nombrosos conceptes tècnics en, per exemple, els models elàstics dels objectes, el procediment del simulador, o en simples formalismes d'expressivitat matemàtica. \newline
	\subsubsection{Problemes temporals}
	Les limitacions temporals són inherents a tot projecte de fi de grau, i aquest no n'és una excepció. Si hi ha una mala gestió del temps, o algun tipus d'imprevist, podem trobar-nos amb la impossibilitat de complir tots els requisits i objectius del treball a temps. \newline
	Per a aquest motiu és important escollir i seguir una metodologia, creant una planificació per a dur a terme el projecte. D'aquesta manera podem ser molt més conscients de l'avenç que s'està fent, i de ser necessari poder realitzar a temps modificacions a l'abast del mateix.
	
	\section{Metodologia i rigor}
	En qualsevol projecte de mínima magnitud, s'ha de triar i adherir-se a un sistema organitzatiu concret. D'aquesta manera podem tenir una visió molt més acurada del treball, i mantenir un ordre en aquest; a més de tenir sempre un procediment a seguir en cas de qualsevol inconvenient.
	\subsection{Metodologia}
	La metodologia emprada en el projecte per a la seva organització i concepció és la \textit{Kanban} \cite{Anderson2010}. Aquesta és una metodologia àgil molt semblant a \textit{Scrum}, però dissenyada per a mantenir un flux constant de resultats evitant discretitzar el projecte en \textit{sprints}. Aquest mètode aporta molta flexibilitat i permet focalitzar el treball en només el que es considera essencial.
	\newline
	La principal idea darrere d'aquest mètode és mantenir sempre alguna tasca en procés, de manera que quan s'ha acabat quelcom, la següent tasca de més prioritat entra en acció. A més a més, s'ha d'evitar mantenir més d'una tasca en procés alhora, limitant la quantitat de \textit{work in progress} al mínim, i centrar-se en el més important en aquell moment.
	\newline
	En cas de tenir qualsevol inconvenient, aquesta metodologia és molt responsiva als canvis i ens permet, ràpidament, corregir qualsevol aspecte mancant, alhora que es valora la seva prioritat amb el director del projecte.
	
	\subsection{Eines de seguiment}
	Per al correcte seguiment del projecte, el seu estat i les seves versions, s'usen diverses eines. Primerament és important tenir una visió global del que es porta fet, i el que queda per fer; per a això s'usa el mètode de \textit{Bullet Journaling} per a ser conscients en tot moment de l'estat del projecte, a un nivell personal. \newline
	A escala de codi, un historial de versions, organitzat adequadament per funcionalitats, és emmagatzemat a \textit{GitHub} juntament amb les diverses versions de la memòria del projecte. Amb aquesta eina podem tenir una visió ordenada del procés evolutiu del treball, veient les funcionalitats per separat.\newline
	Finalment, també es duran a terme reunions periòdiques amb el director on es revisaran, entre altres coses, els avanços fets des de l'última reunió, l'estat del projecte, i dubtes i qüestions sorgides. Donada la naturalesa del projecte, en aquestes reunions es duran a terme explicacions teòriques de conceptes no vistos a la FIB per al correcte desenvolupament del projecte.
	\newline
	En cas de sorgir problemes o imprevistos, es contactarà amb el director i, de ser necessari, es crearà un pla de contingència o una ruta a seguir per a solucionar qualsevol inconvenient.
	 
	
	
	
	% REFERENCIES
	\newpage
	
	\renewcommand{\bibname}{Referències}
	
	\bibliography{BibTex/cit} 
	\bibliographystyle{ieeetr}
	
\end{document}